% ===== PRÉAMBULE CANONIQUE — à coller VERBATIM en tête de CHAQUE .tex =====
\documentclass[11pt,a4paper]{article}
\usepackage[utf8]{inputenc}\usepackage[T1]{fontenc}\usepackage{lmodern}
\usepackage[french]{babel}
\usepackage{amsmath,amssymb,mathtools}\usepackage{icomma}
\usepackage{siunitx}\sisetup{locale=FR}  % \qty{}{}, \si{}, \ang{}
\usepackage{xcolor}\usepackage{colortbl}
\usepackage{tikz}\usepackage{pgfplots}\pgfplotsset{compat=1.18}
\usepackage[siunitx,europeanresistors]{circuitikz} % circuits (élec.)
\usepackage[most]{tcolorbox}\usepackage{enumitem}\usepackage{array,booktabs}
\usepackage[a4paper,margin=2.2cm]{geometry}
\usetikzlibrary{arrows.meta,calc,angles,quotes,positioning,patterns,%
  backgrounds,shapes.geometric,decorations.pathmorphing,decorations.markings,intersections}
\definecolor{primary}{RGB}{222,49,99}\definecolor{primarydark}{RGB}{150,22,55}
\definecolor{defcol}{RGB}{0,114,178}\definecolor{methcol}{RGB}{52,92,148}
\definecolor{excol}{RGB}{0,135,90}\definecolor{attncol}{RGB}{200,40,55}
\tcbset{boxbase/.style={enhanced,breakable,boxrule=0.9pt,arc=2.5pt,
  left=3mm,right=3mm,top=2.3mm,bottom=2.3mm,fonttitle=\bfseries,coltitle=white}}
% --- boîtes TUTORIEL (noms EXACTS, ne pas renommer) ---
\newtcolorbox{cle}[1][]{boxbase,colback=defcol!6,colframe=defcol,
  title={\textsc{À retenir}\ifx\relax#1\relax\relax\else\textmd{ : \itshape #1}\fi}}
\newtcolorbox{methode}[1][]{boxbase,colback=methcol!7,colframe=methcol,sharp corners,
  title={\textsc{Méthode}\ifx\relax#1\relax\relax\else\textmd{ --- \itshape #1}\fi}}
\newtcolorbox{exemplebox}[1][]{boxbase,colback=excol!5,colframe=excol,
  title={\textsc{Exemple}\ifx\relax#1\relax\relax\else\textmd{ : \itshape #1}\fi}}
\newtcolorbox{piege}[1][]{boxbase,colback=attncol!6,colframe=attncol,
  title={\textsc{Piège}\ifx\relax#1\relax\relax\else\textmd{ : \itshape #1}\fi}}
\newtcolorbox{remarque}[1][]{boxbase,colback=black!4,colframe=black!45,
  title={\textsc{Remarque}\ifx\relax#1\relax\relax\else\textmd{ : \itshape #1}\fi}}
% --- boîtes EXERCICE / CORRIGÉ (noms EXACTS, auto-comptées) ---
\newtcolorbox[auto counter]{exo}[1][]{boxbase,colback=primary!4,colframe=primary,
  title={\textsc{Exercice}~\thetcbcounter\ifx\relax#1\relax\relax\else\textmd{ --- \itshape #1}\fi}}
\newtcolorbox[auto counter]{corrige}[1][]{boxbase,colback=excol!4,colframe=excol,
  title={\textsc{Corrigé}~\thetcbcounter\ifx\relax#1\relax\relax\else\textmd{ --- \itshape #1}\fi}}
\newcommand{\vv}[1]{\vec{#1}}\newcommand{\ds}{\displaystyle}
\newcommand{\R}{\mathbb{R}}\newcommand{\N}{\mathbb{N}}\newcommand{\Z}{\mathbb{Z}}
% (un \usepackage{fancyhdr}+en-tête est possible ; il est ignoré côté web)
% =========================================================================
\begin{document}
\begin{corrige}[deux forces, deux configurations]
\begin{enumerate}[label=\alph*)]
  \item \textbf{Correcte.} La résultante de deux forces parallèles est toujours parallèle à
        celles-ci, que les forces soient de même sens ou de sens opposé.
  \item \textbf{Correcte.} Cas b) (sens opposé) : $F_r = F_2 - F_1 = 2{,}5 - 1 = \qty{1.5}{N}$.
  \item \textbf{Correcte.} Cas a) : $\dfrac{\overline{AO}}{\overline{BO}}=\dfrac{F_2}{F_1}=2{,}5$ et
        $\overline{AO}+\overline{BO}=70$, d'où $3{,}5\,\overline{BO}=70$, soit $\overline{BO}=\qty{20}{cm}$.
        $O$ est bien entre $A$ et $B$, à \qty{20}{cm} de $B$.
  \item \textbf{Incorrecte.} En sens opposé, $O$ sort du segment \emph{du côté de la plus grande
        force} $F_2$, c'est-à-dire du côté de $B$ — et non du côté de $A$. \emph{C'est la proposition
        fausse.}
\end{enumerate}
\end{corrige}

\begin{corrige}[tige maintenue horizontale par une masse]
À l'équilibre de translation, le poids équilibre les deux forces montantes :
$G = F_1+F_2 = 1+2{,}5 = \qty{3.5}{N}$.
\begin{enumerate}[label=\alph*)]
  \item \textbf{Fausse.} $m = \dfrac{G}{g} = \dfrac{3{,}5}{10} = \qty{0.35}{kg} = \qty{350}{g}$
        (et non \qty{35}{g}).
  \item \textbf{Correcte.} $\dfrac{\overline{AO}}{\overline{BO}}=\dfrac{F_2}{F_1}=\dfrac{2{,}5}{1}=2{,}5$.
  \item \textbf{Correcte.} $O$ est le point d'application de la résultante des deux forces : la relation
        $F_1\times\overline{AO}=F_2\times\overline{BO}$ signifie exactement que les moments de $\vv F_1$
        et $\vv F_2$ autour de $O$ se compensent ; de plus $\vv G$ passe par $O$ (bras nul). La somme
        des moments autour de $O$ est donc nulle.
  \item \textbf{Fausse.} Puisque $F_1\times\overline{AO}=F_2\times\overline{BO}$, on a
        $F_2\times\overline{BO}-F_1\times\overline{AO}=0$, et $G\times 0 = 0$ : la somme vaut
        \textbf{exactement 0}, elle n'est pas strictement positive.
\end{enumerate}
\end{corrige}

\begin{corrige}[composition de trois forces parallèles]
\begin{enumerate}[label=\alph*)]
  \item Même sens : $R_{12}=F_1+F_2=10+20=\qty{30}{N}$. Position : $\overline{AO_1}/\overline{BO_1}
        =F_2/F_1=2$ et $\overline{AO_1}+\overline{BO_1}=30$, d'où $\overline{BO_1}=\qty{10}{cm}$ et
        $\overline{AO_1}=\qty{20}{cm}$ : $\vv R_{12}$ s'applique à $x=\qty{20}{cm}$.
  \item On compose $\vv R_{12}$ ($\qty{30}{N}$ à $x=20$) avec $\vv F_3$ ($\qty{30}{N}$ à $x=60$). Les
        deux intensités sont \emph{égales} : le point d'application est au \textbf{milieu}, à
        $x=\dfrac{20+60}{2}=\qty{40}{cm}$. Résultante globale :
        \[ R = R_{12}+F_3 = 30+30 = \qty{60}{N}\ \text{à}\ x=\qty{40}{cm}. \]
\end{enumerate}
\end{corrige}

\begin{corrige}[résultante et force équilibrante]
\begin{enumerate}[label=\alph*)]
  \item Sens opposé : $F_r = F_2 - F_1 = 8 - 3 = \qty{5}{N}$, de même sens que $\vv F_2$ (vers le haut).
  \item $O$ est hors de $[AB]$, du côté de $B$ (force la plus grande). En posant $\overline{BO}=x$ :
        $\overline{AO}=40+x$ et $F_1\times\overline{AO}=F_2\times\overline{BO}$ :
        \[ 3(40+x)=8x \Rightarrow 120=5x \Rightarrow x=\overline{BO}=\qty{24}{cm}. \]
        $O$ est à \qty{24}{cm} au-delà de $B$.
  \item La force équilibrante s'applique en $O$, de même intensité que $\vv F_r$ mais de sens opposé :
        \qty{5}{N} \textbf{vers le bas}. Vérification : forces vers le haut $= F_2 = \qty{8}{N}$ ;
        vers le bas $= F_1 + F_{\text{éq}} = 3 + 5 = \qty{8}{N}$. Elles se compensent :
        $\sum\vv F=\vv 0$.
\end{enumerate}
\end{corrige}

\end{document}
